\title[La gracia]{\bfseries{La gracia de Dios}}
\subtitle{``Pero Noé halló gracia ante los ojos de \SENOR{}'' (Gn 6:8)}
\author{}\institute{}\date{}
\titlegraphic{\includegraphics[width=0.80\linewidth,height=0.36\textheight,keepaspectratio]{portada.jpg}\\[-1pt]
  {\fontsize{5}{6}\selectfont\color{gris}\textit{Nota.} Elaboración propia: Gn 6:8 en texto hebreo consonántico (BHS).}}
\setbeamerfont{section in toc}{size=\small}
\AtBeginSection[]{\begin{frame}\frametitle{Contenidos}\tableofcontents[currentsection]\end{frame}}
\begin{document}
\frame{\titlepage}

\begin{frame}\frametitle{Contenidos}\tableofcontents\end{frame}

% ------------------------------------------------------------------
\section{Introducción}

\begin{frame}{Un favor que no se debe}
\begin{columns}[c,onlytextwidth]
\begin{column}{0.50\textwidth}
\begin{tcolorbox}[colback=papiro,colframe=dorado,boxrule=1.2pt,arc=3pt]
\centering\large\itshape Favor inmerecido de Dios hacia pecadores.
\end{tcolorbox}
\small\RaggedRight Tres partes: \destacar{favor} (Dios se inclina y da), \destacar{inmerecido} (no paga una deuda) y \destacar{hacia pecadores} (alcanza a quien merecía lo contrario).
\end{column}
\begin{column}{0.48\textwidth}
\scriptsize
\begin{cita}Cuando se dice que una cosa es por ``gracia'', queremos decir que el receptor de la misma no tiene derecho a reclamarla, que de ninguna manera era una deuda.\fuente{(Pink, 2020, cap. 13, p. 68)}\end{cita}
\begin{cita}Es el eterno y absoluto favor gratuito de Dios, manifestado en conceder bendiciones espirituales y eternas a los culpables e indignos.\fuente{(A. Booth, en Pink, 2020, p. 68)}\end{cita}
\precision{la definición del recuadro es la de la serie; el desarrollo es una propuesta sobre fuentes verificadas.}
\end{column}
\end{columns}
\end{frame}

\begin{frame}{Gloriarse solo en Dios}
\begin{cita}Mas alábese en esto el que se hubiere de alabar: en entenderme y conocerme, que yo soy \SENOR{}\ldots\fuente{(Jer 9:24 [9:23], RVR1960)}\end{cita}
\begin{cita}¿qué tienes que no hayas recibido? Y si lo recibiste, ¿por qué te glorías como si no lo hubieras recibido?\fuente{(1 Co 4:7, RVR1960)}\end{cita}
\small La gracia es el atributo que más directamente quita al ser humano \destacar{todo otro motivo de gloria}: ``no por obras, para que nadie se gloríe'' (Ef 2:9).\par\medskip
\centering
\begin{tikzpicture}
 \foreach \x/\n/\t/\r in {0/1/{La palabra}/{Gn 6:8; Éx 34:6}, 3.15/2/{Cantada y narrada}/{Sal 51; 2 Cr 33}, 6.3/3/{En Cristo}/{Ro 3:24; 1 Co 15:10}, 9.45/4/{Vivirla}/{Mt 10:8; He 4:16}}{
  \node[caja,anchor=north] (b\n) at (\x,0) {\parbox[t][0.9cm][t]{2.55cm}{\RaggedRight
    {\small\textbf{\textcolor{azul}{\t}}}\par\vfill{\tiny\color{gris}(\r)}}};
  \node[circle,fill=azul,text=papiro,font=\scriptsize\bfseries,inner sep=1.6pt] at (b\n.north west) {\n};}
\end{tikzpicture}
\end{frame}

% ------------------------------------------------------------------
\section{Génesis 6}

\begin{frame}{Texto ancla: Gn 6:5--8}
\fotolado{panel_genesis6.jpg}{Elaboración propia: Gn 6:8 en texto hebreo consonántico (BHS).}
\begin{columns}[c,onlytextwidth]
\begin{column}{0.58\textwidth}
\begin{cita}Y vio \SENOR{} que la maldad de los hombres era mucha en la tierra, y que todo designio de los pensamientos del corazón de ellos era de continuo solamente el mal\ldots{} Y dijo \SENOR: Raeré de sobre la faz de la tierra a los hombres que he creado\ldots\fuente{(Gn 6:5, 7, RVR1960)}\end{cita}
\begin{cita}Pero Noé halló gracia ante los ojos de \SENOR{}.\fuente{(Gn 6:8, RVR1960)}\end{cita}
\small La \destacar{primera} vez que la Biblia nombra la gracia es en el \destacar{peor momento} de la humanidad.
\end{column}
\begin{column}{0.40\textwidth}\end{column}
\end{columns}
\end{frame}

\begin{frame}{Gn 6:8 en hebreo}
\begin{center}
{\Large\heb{וְנֹחַ מָצָא חֵן בְּעֵינֵי יְהוָה}}
\end{center}
\renewcommand{\arraystretch}{1.15}\scriptsize
\begin{tabular}{@{}>{\RaggedRight\arraybackslash}p{2.9cm}>{\RaggedRight\arraybackslash}p{1.3cm}>{\RaggedRight\arraybackslash}p{7.6cm}@{}}
\toprule
\textbf{Término} & \textbf{Raíz} & \textbf{Sentido} \\ \midrule
\heb{וְנֹחַ} \emph{veNóaj} & \heb{נוח} & ``Pero Noé'': el nombre antes del verbo marca contraste con 6:5--7. \\
\heb{מָצָא חֵן} \emph{matsá jen} & \heb{מצא} & ``halló gracia'': fórmula del que recibe el favor de un superior. No la produjo: la \destacar{encontró}. \\
\heb{חֵן} \emph{jen} & \heb{חנן} & ``gracia'', ``favor'': primera aparición en la Biblia. \\
\heb{בְּעֵינֵי יְהוָה} & \heb{עין} & ``en los ojos del \SENOR''. La mirada que ``vio'' la maldad (6:5) se posa con favor. \\
\grc{Νωε δὲ εὗρεν χάριν} & \grc{χάρις} & la Septuaginta traduce \emph{jen} con \grc{χάρις}, la palabra de Pablo. \\
\bottomrule
\end{tabular}
\end{frame}

\begin{frame}{\emph{Nóaj} y \emph{jen}: un espejo de letras}
\centering
\begin{tikzpicture}[font=\small]
 \node[draw=azul,line width=1pt,fill=papiro,rounded corners=3pt,minimum width=3.4cm,minimum height=1.9cm,align=center] (n) at (0,0) {\Huge\heb{נח}\\[3pt]\small \emph{Nóaj}, ``Noé''};
 \node[draw=dorado,line width=1.6pt,fill=azul,text=papiro,rounded corners=3pt,minimum width=3.4cm,minimum height=1.9cm,align=center] (j) at (6.4,0) {\Huge\heb{חן}\\[3pt]\small \emph{jen}, ``gracia''};
 \draw[{Latex[length=3mm]}-{Latex[length=3mm]},dorado,line width=1.4pt] (n.east) -- node[above,font=\scriptsize,text=gris] {mismas consonantes} node[below,font=\scriptsize,text=gris] {orden inverso} (j.west);
\end{tikzpicture}\par\medskip
\small\RaggedRight En el rollo, sin vocales, el nombre y la palabra se miran como en un espejo: un juego \destacar{visual} más que sonoro (Sasson, 1975; cf. Cassuto, 1961).\par\smallskip
\precision{el juego de letras es observación de los comentaristas, no afirmación del texto. Lo que el texto afirma es el contraste: toda la humanidad bajo sentencia y un hombre que ``halló gracia''.}
\end{frame}

\begin{frame}{La gracia antes de la justicia}
\centering
\begin{tikzpicture}[node distance=0.35cm]
 \node[caja,text width=2.45cm,align=center,minimum height=1.35cm] (a) {\textbf{6:5--7}\\ toda la humanidad bajo sentencia};
 \node[caja,draw=dorado,line width=1.4pt,text width=2.45cm,align=center,minimum height=1.35cm,right=of a] (b) {\textbf{6:8}\\ ``Noé \destacar{halló gracia}''};
 \node[caja,text width=2.45cm,align=center,minimum height=1.35cm,right=of b] (c) {\textbf{6:9}\\ ``Noé, varón justo\ldots{} con Dios caminó''};
 \node[caja,text width=2.45cm,align=center,minimum height=1.35cm,right=of c] (d) {\textbf{8:21}\\ ``el corazón del hombre es malo'': marco de promesa};
 \draw[-{Latex},dorado,line width=1.2pt] (a)--(b); \draw[-{Latex},dorado,line width=1.2pt] (b)--(c); \draw[-{Latex},dorado,line width=1.2pt] (c)--(d);
\end{tikzpicture}\par\medskip
\small\RaggedRight
\begin{itemize}
\item Primero la gracia (6:8); después, en una nueva sección, la justicia de Noé (6:9).
\item La humanidad \destacar{no cambió} después del diluvio (8:21): Dios decidió tratarla con favor.
\end{itemize}
\precision{que el orden 6:8--9 indique que la gracia precede a la justicia es lectura teológica: 6:9 abre una nueva sección (\emph{toledot}). La tradición judía discutió la pregunta (ver ``Judaísmo'').}
\end{frame}

% ------------------------------------------------------------------
\section{La palabra}

\begin{frame}{La raíz \heb{חנן} (\emph{janán})}
\begin{columns}[t,onlytextwidth]
\begin{column}{0.58\textwidth}
\scriptsize\renewcommand{\arraystretch}{1.1}
\begin{tabular}{@{}>{\RaggedRight\arraybackslash}p{2.3cm}>{\RaggedRight\arraybackslash}p{2.9cm}>{\RaggedRight\arraybackslash}p{1.5cm}@{}}
\toprule
\textbf{Forma} & \textbf{Uso} & \textbf{Texto} \\ \midrule
\heb{שֶׁקֶר הַחֵן} & ``engañosa es la gracia'' & Pr 31:30 \\
\heb{יַעֲלַת־חֵן} & ``graciosa gacela'' & Pr 5:19 \\
\heb{מָצָא חֵן בְּעֵינֵי} & favor de un superior & Gn 39:4 \\
\heb{אֲשֶׁר־חָנַן אֱלֹהִים} & ``los niños que Dios ha dado'' & Gn 33:5 \\
\heb{חָנֵּנִי} & ``ten piedad de mí'' & Sal 51:1 \\
\heb{וָאֶתְחַנַּן} & ``y supliqué gracia'' & Dt 3:23 \\
\heb{תְּחִנָּתוֹ} & ``su súplica'' & 2 Cr 33:13 \\
\heb{חִנָּם} & ``de balde'', ``gratis'' & Éx 21:2 \\
\bottomrule
\end{tabular}
\end{column}
\begin{column}{0.40\textwidth}\small\RaggedRight
Tres planos:
\begin{itemize}
\item \destacar{Lo que se ve}: el encanto que agrada.
\item \destacar{El favor}: siempre se halla ``en los ojos'' del que lo da.
\item \destacar{La súplica}: el que no puede reclamar, pide gracia.
\end{itemize}
La gracia nace en Dios, desciende desde arriba y la recibe quien solo puede pedirla.
\end{column}
\end{columns}
\end{frame}

\begin{frame}{``¿Por qué he hallado gracia?''}
\begin{columns}[c,onlytextwidth]
\begin{column}{0.54\textwidth}
\centering\heb{מַדּוּעַ מָצָאתִי חֵן בְּעֵינֶיךָ לְהַכִּירֵנִי וְאָנֹכִי נָכְרִיָּה}
\begin{cita}¿Por qué he hallado gracia en tus ojos para que me reconozcas, siendo yo extranjera?\fuente{(Rt 2:10, RVR1960)}\end{cita}
\end{column}
\begin{column}{0.44\textwidth}\small\RaggedRight
\begin{itemize}
\item \heb{לְהַכִּירֵנִי} \emph{lehakirení}, ``que me reconozcas'', y \heb{נָכְרִיָּה} \emph{nojriyyá}, ``extranjera'': la \destacar{misma raíz} \heb{נכר}.
\item La que no tiene derecho a ser reconocida, es reconocida.
\item Rut no pregunta qué hizo; pregunta \destacar{por qué}. Esa pregunta es la marca de la gracia.
\end{itemize}
\end{column}
\end{columns}
\end{frame}

\begin{frame}{\heb{חִנָּם} (\emph{jinnam}): de balde}
\scriptsize\renewcommand{\arraystretch}{1.15}
\begin{tabular}{@{}>{\RaggedRight\arraybackslash}p{2.0cm}>{\RaggedRight\arraybackslash}p{7.2cm}>{\RaggedRight\arraybackslash}p{2.4cm}@{}}
\toprule
\textbf{Sentido} & \textbf{Ejemplo (RVR1960)} & \textbf{Texto} \\ \midrule
Sin pago & ``¿me servirás de balde? Dime cuál será tu salario'' & Gn 29:15 \\
Sin pago & ``al séptimo saldrá libre, de balde'' & Éx 21:2 \\
Sin costo & ``holocaustos que no me cuesten nada'' & 2 S 24:24 \\
Sin motivo & ``¿Acaso teme Job a Dios de balde?'' & Job 1:9 \\
Sin causa & ``los que me aborrecen sin causa'' & Sal 69:4 [69:5] \\
Sin precio & ``De balde fuisteis vendidos; por tanto, sin dinero seréis rescatados'' & Is 52:3 \\
\bottomrule
\end{tabular}\par\medskip
\small Adverbio formado de \emph{jen}: ``por gracia'' = \destacar{gratis}. Jesús aplica el Sal 69 a sí mismo: ``Sin causa (\grc{δωρεάν}) me aborrecieron'' (Jn 15:25). Pablo: ``justificados \destacar{gratuitamente} (\grc{δωρεάν})'' (Ro 3:24).\par\smallskip
\precision{la unión teológica de ambos sentidos (aborrecido de balde, justificados de balde) es lectura propia; el vínculo léxico es la LXX.}
\end{frame}

\begin{frame}{El espectro de ``gracia''}
\scriptsize\renewcommand{\arraystretch}{1.08}
\begin{tabular}{@{}>{\RaggedRight\arraybackslash}p{2.7cm}>{\RaggedRight\arraybackslash}p{1.0cm}>{\RaggedRight\arraybackslash}p{5.0cm}>{\RaggedRight\arraybackslash}p{2.8cm}@{}}
\toprule
\textbf{Término} & \textbf{Veces} & \textbf{En la RVR1960} & \textbf{Textos} \\ \midrule
\heb{חֵן} \emph{jen} & $\approx$\,69 & ``gracia'' (casi siempre), ``favor'' & Gn 6:8; Zac 4:7 \\
\heb{חָנַן} \emph{janán} & $\approx$\,78 & ``tener misericordia'', ``piedad''; ``rogar'' & Éx 33:19; Nm 6:25 \\
\heb{חַנּוּן} \emph{jannun} & 13 & ``clemente'', ``misericordioso'', ``piadoso'' & Éx 22:27; 34:6 \\
\heb{תְּחִנָּה}; \heb{תַּחֲנוּנִים} & 25; 18 & ``ruego'', ``súplica'', ``oración'' & 1 R 8:28; Zac 12:10 \\
\heb{חִנָּם} \emph{jinnam} & 32 & ``de balde'', ``sin causa'' & Éx 21:2; Is 52:3 \\ \midrule
\grc{χάρις} \emph{járis} & 155 & ``gracia'', ``gracias'', ``favor'', ``mérito'' & Ro 3:24; Lc 6:32 \\
\grc{χάρισμα}; \grc{χαρίζομαι} & 17; 23 & ``don'', ``dádiva''; ``perdonar'', ``conceder'' & Ro 6:23; Ef 4:32 \\
\grc{δωρεάν} \emph{dōreán} & 9 & ``gratuitamente'', ``de gracia'', ``por demás'' & Ro 3:24; Gá 2:21 \\
\bottomrule
\end{tabular}\par\smallskip
\small ``Gracia'' en la RVR1960: \destacar{211 veces} en 199 vv (74 AT, 137 NT). El lector español \destacar{no ve} la gracia en Éx 33:19, Nm 6:25 o Sal 51:1: la RVR1960 traduce \emph{janán} ``misericordia''.\par
\precision{conteo propio (RVR1960 sin títulos, BHS y NA27 con Strong); las cifras $\approx$ son las de las concordancias. Verificar en e-Sword.}
\end{frame}

\begin{frame}{Del hebreo a Pablo}
\centering
\begin{tikzpicture}[font=\small,node distance=0.75cm]
 \node[caja,text width=3.3cm,align=center,font=\small] (h) {\textbf{Hebreo}\\ \heb{חֵן} (Gn 6:8)\\ \heb{חִנָּם} (Is 52:3)\\ \heb{חָנַן}, \heb{חַנּוּן} (Éx 34:6)};
 \node[caja,text width=3.3cm,align=center,font=\small,right=of h] (l) {\textbf{Septuaginta}\\ \grc{χάρις}\\ \grc{δωρεάν}\\ \grc{ἐλεέω}, \grc{ἐλεήμων}};
 \node[draw=dorado,line width=1.6pt,fill=azul,text=papiro,rounded corners=3pt,text width=3.3cm,align=center,inner sep=5pt,right=of l] (p) {\textbf{Pablo}\\ \grc{χάρις} (Ro 3:24)\\ \grc{δωρεάν} (Ro 3:24)\\ \grc{ἐλεέω} (Ro 9:15)};
 \draw[-{Latex[length=3mm]},dorado,line width=1.4pt] (h)--(l);
 \draw[-{Latex[length=3mm]},dorado,line width=1.4pt] (l)--(p);
\end{tikzpicture}\par\medskip
{\small\itshape\color{gris} La LXX reparte la raíz: \emph{jen} $\to$ \grc{χάρις}; \emph{janán}, \emph{jannun} $\to$ familia de \grc{ἔλεος}.}\par\medskip
\small\RaggedRight En el griego común, \grc{χάρις} era el favor que obliga a devolver. Pablo la da a \destacar{quien no tiene con qué devolverla}, ``a los impíos'' (Ro 5:6), en la línea de \emph{jen} y \emph{jinnam}.\par\smallskip
\precision{lectura de Barclay (2015), que subraya la ``incongruencia'' del don en Pablo; otros acentúan la continuidad con la \emph{jésed} del pacto (Dunn, 1998).}
\end{frame}

% ------------------------------------------------------------------
\section{El Nombre}

\begin{frame}{``Porque soy clemente''}
\begin{columns}[c,onlytextwidth]
\begin{column}{0.56\textwidth}
\centering\heb{וְהָיָה כִּי־יִצְעַק אֵלַי וְשָׁמַעְתִּי כִּי־חַנּוּן אָנִי}
\begin{cita}Si tomares en prenda el vestido de tu prójimo, a la puesta del sol se lo devolverás. Porque sólo eso es su cubierta\ldots{} ¿En qué dormirá? Y cuando él clamare a mí, yo le oiré, porque soy misericordioso.\fuente{(Éx 22:26--27 [22:25--26], RVR1960)}\end{cita}
\end{column}
\begin{column}{0.42\textwidth}\small\RaggedRight
\begin{itemize}
\item \heb{כִּי־חַנּוּן אָנִי}, ``porque \destacar{clemente} soy yo'': el primer \emph{jannun} de la Biblia, dicho por Dios de sí mismo.
\item El pobre \destacar{no tiene derecho} a reclamar: la prenda es legal.
\item Antes del Sinaí, Dios se presenta como el que se inclina hacia quien no tiene nada que ofrecer.
\end{itemize}
\end{column}
\end{columns}
\end{frame}

\begin{frame}{\emph{El rajum vejannun}}
\fotolado{panel_exodo34.jpg}{Elaboración propia: Éx 34:6 en texto hebreo consonántico (BHS).}
\begin{columns}[c,onlytextwidth]
\begin{column}{0.58\textwidth}
\begin{cita}tendré misericordia del que tendré misericordia, y seré clemente para con el que seré clemente.\fuente{(Éx 33:19, RVR1960)}\end{cita}
\begin{cita}Y pasando \SENOR{} por delante de él, proclamó: ¡\SENOR! ¡\SENOR! fuerte, misericordioso y piadoso; tardo para la ira, y grande en misericordia y verdad.\fuente{(Éx 34:6, RVR1960)}\end{cita}
\scriptsize \heb{אֵל רַחוּם וְחַנּוּן}: ``Dios compasivo y \destacar{clemente}''. Proclamado justo \destacar{después del becerro}, ante un pueblo ``de dura cerviz'' (34:9).
\end{column}
\begin{column}{0.40\textwidth}\end{column}
\end{columns}
\end{frame}

\begin{frame}{``Tendré gracia de quien tendré gracia''}
\begin{columns}[t,onlytextwidth]
\begin{column}{0.52\textwidth}
\centering\heb{וְחַנֹּתִי אֶת־אֲשֶׁר אָחֹן וְרִחַמְתִּי אֶת־אֲשֶׁר אֲרַחֵם}\par\medskip
\RaggedRight\small
\begin{itemize}
\item Construcción \emph{idem per idem}, como ``Yo soy el que soy'' (Éx 3:14): la causa de la gracia está \destacar{en Dios}.
\item Moisés apela a ella sin presentar mérito: ``vaya ahora el Señor en medio de nosotros; \destacar{porque} es un pueblo de dura cerviz'' (Éx 34:9).
\end{itemize}
\end{column}
\begin{column}{0.46\textwidth}
\begin{cita}Y oró a \SENOR{} y dijo\ldots{} sabía yo que tú eres Dios clemente y piadoso, tardo en enojarte, y de grande misericordia, y que te arrepientes del mal.\fuente{(Jon 4:2, RVR1960)}\end{cita}
\scriptsize Jonás conoce el Nombre y \destacar{por eso huye}: la gracia escandaliza a quien quiere que el otro reciba lo que merece.
\end{column}
\end{columns}
\end{frame}

\begin{frame}{Los nombres que revelan la gracia}
\scriptsize\renewcommand{\arraystretch}{1.18}
\begin{tabular}{@{}>{\RaggedRight\arraybackslash}p{2.7cm}>{\RaggedRight\arraybackslash}p{3.4cm}>{\RaggedRight\arraybackslash}p{1.9cm}>{\RaggedRight\arraybackslash}p{3.6cm}@{}}
\toprule
\textbf{Nombre o título} & \textbf{Hebreo / griego} & \textbf{Texto} & \textbf{Qué revela} \\ \midrule
Clemente & \heb{כִּי־חַנּוּן אָנִי} & Éx 22:27 & oye al pobre que no puede reclamar \\
Dios misericordioso y clemente & \heb{אֵל רַחוּם וְחַנּוּן} & Éx 34:6 & su Nombre ante un pueblo infiel \\
Dios clemente y piadoso & \heb{אֵל־חַנּוּן וְרַחוּם} & Jon 4:2 & perdona incluso al enemigo \\
Espíritu de gracia & \heb{רוּחַ חֵן וְתַחֲנוּנִים} & Zac 12:10; He 10:29 & da gracia y mueve a la súplica \\
Trono de la gracia & \grc{ὁ θρόνος τῆς χάριτος} & He 4:16 & el tribunal convertido en socorro \\
Dios de toda gracia & \grc{ὁ θεὸς πάσης χάριτος} & 1 P 5:10 & llama a su gloria \\
\bottomrule
\end{tabular}\par\medskip
\precision{\emph{Shemot Rabá} 3:6: ``según mis obras soy llamado''. ``Clemente'' nace de una obra: oír al que clama sin derecho.}
\end{frame}

% ------------------------------------------------------------------
\section{Definición}

\begin{frame}{Definición}
\begin{tcolorbox}[colback=papiro,colframe=dorado,boxrule=1.2pt,arc=3pt]
\centering\itshape Favor inmerecido de Dios hacia pecadores.
\end{tcolorbox}
\begin{tcolorbox}[colback=papiro,colframe=azul,boxrule=1pt,arc=3pt,fonttitle=\bfseries\small,title=Desarrollada,fontupper=\small]
La gracia de Dios es la perfección eterna por la cual Dios, libre y soberanamente, concede su favor a quienes no lo merecen y merecían lo contrario. No responde a mérito alguno ni crea una deuda: da perdón, vida y comunión ``de balde'', y su fuente y su cauce es Jesucristo (Gn 6:8; Éx 34:6; Ro 3:24; 5:15; Ef 2:8--9).
\end{tcolorbox}
\scriptsize\textbf{Clasificación:} Relación. \emph{Incomunicable} como fuente (solo Dios justifica al impío, Ro 4:5); \emph{comunicable} como reflejo (Mt 10:8; Ef 4:32; Col 4:6).\par
\precision{la definición desarrollada y la clasificación son una propuesta; conviene confirmarlas.}
\end{frame}

\begin{frame}{Rasgos de la gracia}
\scriptsize\renewcommand{\arraystretch}{1.15}
\begin{tabular}{@{}>{\RaggedRight\arraybackslash}p{2.6cm}>{\RaggedRight\arraybackslash}p{5.8cm}>{\RaggedRight\arraybackslash}p{3.2cm}@{}}
\toprule
\textbf{Rasgo} & \textbf{Qué afirma} & \textbf{Textos} \\ \midrule
Favor & inclinación de Dios que se traduce en dones & Gn 6:8; Rt 2:10 \\
Inmerecida & no responde a obras ni a méritos & Dt 9:5--6; Ro 11:6 \\
Hacia pecadores & alcanza a quien merecía el juicio & Gn 6:5--8; Ro 5:6--8 \\
Gratuita & ``de balde'': no se compra ni se devuelve & Is 52:3; 55:1; Ro 3:24 \\
Soberana & Dios la da a quien quiere & Éx 33:19; Ro 9:15--16 \\
Eterna & ``antes de los tiempos de los siglos'' & 2 Ti 1:9 \\
Abundante y reinante & sobreabunda donde abundó el pecado & Ro 5:20--21 \\
Eficaz & transforma y enseña a vivir & 1 Co 15:10; Tit 2:11--12 \\
\bottomrule
\end{tabular}\par\medskip
\small Pink: \destacar{eterna, gratuita y soberana}. ``Si la gracia `reina', entonces está en el trono y el ocupante del trono es soberano'' (2020, cap. 13, p. 69).
\end{frame}

\begin{frame}{Gracia y misericordia}
\begin{columns}[t,onlytextwidth]
\begin{column}{0.58\textwidth}
\scriptsize\renewcommand{\arraystretch}{1.18}
\begin{tabular}{@{}>{\RaggedRight\arraybackslash}p{1.7cm}>{\RaggedRight\arraybackslash}p{1.6cm}>{\RaggedRight\arraybackslash}p{3.4cm}@{}}
\toprule
\textbf{Perfección} & \textbf{Hebreo} & \textbf{Mira al ser humano como\ldots} \\ \midrule
Misericordia & \heb{רַחוּם} & miserable: necesita alivio \\
Gracia & \heb{חַנּוּן} & culpable e indigno: recibe lo que no merece \\
Paciencia & \heb{אֶרֶךְ אַפַּיִם} & rebelde aún no juzgado \\
Bondad & \heb{רַב־חֶסֶד} & criatura que recibe el bien \\
\bottomrule
\end{tabular}
\end{column}
\begin{column}{0.40\textwidth}\small\RaggedRight
Formulación común:\par
\destacar{la misericordia no nos da lo que merecemos; la gracia nos da lo que no merecemos.}\par\medskip
\precision{Pink: la gracia ``es ejercida sólo sobre los elegidos''; otros reformados hablan además de ``gracia común'' (Berkhof, 1938). La serie no toma partido. En hebreo \emph{rajum} y \emph{jannun} van juntos.}
\end{column}
\end{columns}
\end{frame}

% ------------------------------------------------------------------
\section{De balde}

\begin{frame}{El esclavo que sale libre}
\fotolado{panel_isaias52.jpg}{Elaboración propia: Is 52:3 en texto hebreo consonántico (BHS).}
\begin{columns}[c,onlytextwidth]
\begin{column}{0.58\textwidth}
\begin{cita}Si comprares siervo hebreo, seis años servirá; mas al séptimo saldrá libre, de balde.\fuente{(Éx 21:2, RVR1960)}\end{cita}
\begin{cita}Porque así dice \SENOR: De balde fuisteis vendidos; por tanto, sin dinero seréis rescatados.\fuente{(Is 52:3, RVR1960)}\end{cita}
\scriptsize Sale \heb{חִנָּם}: sin pagar rescate. Y no ``con las manos vacías'', ``acordándote de que fuiste siervo\ldots{} y \SENOR{} tu Dios te rescató'' (Dt 15:13--15). LXX Is 52:3: \grc{δωρεὰν ἐπράθητε}.
\end{column}
\begin{column}{0.40\textwidth}\end{column}
\end{columns}
\end{frame}

\begin{frame}{``Comprad sin dinero''}
\begin{columns}[c,onlytextwidth]
\begin{column}{0.54\textwidth}
\centering\heb{לְכוּ שִׁבְרוּ בְּלוֹא־כֶסֶף וּבְלוֹא מְחִיר יַיִן וְחָלָב}
\begin{cita}A todos los sedientos: Venid a las aguas; y los que no tienen dinero, venid, comprad y comed. Venid, comprad sin dinero y sin precio, vino y leche.\fuente{(Is 55:1, RVR1960)}\end{cita}
\end{column}
\begin{column}{0.44\textwidth}\small\RaggedRight
\begin{itemize}
\item \heb{שָׁבַר}: comprar grano en tiempo de hambre (Gn 42:2--3). Comida real, \destacar{sin precio}.
\item El rescate ``sin dinero'' no es gratis para el Siervo: ``herido fue por nuestras rebeliones'' (Is 53:5; 1 P 1:18--19).
\item ``tome del agua de la vida \destacar{gratuitamente}'' (\grc{δωρεάν}, Ap 22:17).
\end{itemize}
\end{column}
\end{columns}
\end{frame}

\begin{frame}{El salario y la dádiva}
\centering
\begin{tikzpicture}[font=\small]
 \begin{scope}
  \node[font=\bfseries,text=gris] at (2.6,3.2) {El salario};
  \fill[papiro] (0,0) rectangle (5.2,2.8);
  \draw[gris,line width=0.8pt] (0,0) rectangle (5.2,2.8);
  \node[text width=4.8cm,align=center,font=\scriptsize] at (2.6,1.4) {\grc{μισθός}, \grc{ὀψώνια}\\[3pt] se cuenta ``como deuda'' (Ro 4:4)\\[3pt] ``la paga del pecado es muerte''\\[3pt] (Gn 29:15; Ro 6:23)};
 \end{scope}
 \begin{scope}[xshift=6.2cm]
  \node[font=\bfseries,text=azul] at (2.6,3.2) {La dádiva};
  \fill[azul!8] (0,0) rectangle (5.2,2.8);
  \draw[azul,line width=1.2pt] (0,0) rectangle (5.2,2.8);
  \node[text width=4.8cm,align=center,font=\scriptsize] at (2.6,1.4) {\grc{χάρις}, \grc{χάρισμα}\\[3pt] se cuenta ``como gracia'' (Ro 4:4)\\[3pt] ``la dádiva de Dios es vida eterna''\\[3pt] (Éx 21:2; Is 55:1; Ro 6:23)};
 \end{scope}
 \draw[-{Latex[length=3mm]},dorado,line width=1.6pt] (5.3,1.4) -- (6.1,1.4);
\end{tikzpicture}\par\medskip
\small\RaggedRight \grc{ὀψώνια}: la paga del soldado (Moo, 1996). El pecado paga puntualmente. Dios no paga: \destacar{regala}. Por eso ``justifica al impío'' (Ro 4:5), frase que en la ley describía a un juez injusto (Éx 23:7; Pr 17:15).
\end{frame}

% ------------------------------------------------------------------
\section{Cantada}

\begin{frame}{``Ten piedad de mí'': Sal 51}
\begin{columns}[c,onlytextwidth]
\begin{column}{0.52\textwidth}
\centering\heb{חָנֵּנִי אֱלֹהִים כְּחַסְדֶּךָ כְּרֹב רַחֲמֶיךָ מְחֵה פְשָׁעָי}
\begin{cita}Ten piedad de mí, oh Dios, conforme a tu misericordia; Conforme a la multitud de tus piedades borra mis rebeliones.\fuente{(Sal 51:1 [51:3], RVR1960)}\end{cita}
\end{column}
\begin{column}{0.46\textwidth}\small\RaggedRight
\begin{itemize}
\item \heb{חָנֵּנִי} \emph{jonnení}: ``ten \destacar{gracia} de mí''.
\item Las tres palabras de Éx 34:6 (\emph{jannun}, \emph{jésed}, \emph{rajamim}) en la primera línea.
\item La medida no es David: ``conforme a \destacar{tu} misericordia''.
\item ``Ten misericordia de nosotros, oh \SENOR, ten misericordia de nosotros'' (\heb{חָנֵּנוּ יְהוָה חָנֵּנוּ}, Sal 123:3).
\end{itemize}
\end{column}
\end{columns}
\end{frame}

\begin{frame}{Gracia y gloria; ``gracia, gracia a ella''}
\begin{columns}[t,onlytextwidth]
\begin{column}{0.49\textwidth}
\centering\heb{חֵן וְכָבוֹד יִתֵּן יְהוָה}
\begin{cita}Porque sol y escudo es \SENOR{} Dios; Gracia y gloria dará \SENOR. No quitará el bien a los que andan en integridad.\fuente{(Sal 84:11 [84:12], RVR1960)}\end{cita}
\scriptsize La gracia es el camino; la gloria, el destino (Ro 5:2; 1 P 5:10).
\end{column}
\begin{column}{0.49\textwidth}
\centering\heb{תְּשֻׁאוֹת חֵן חֵן לָהּ}
\begin{cita}\ldots{} él sacará la primera piedra con aclamaciones de: Gracia, gracia a ella.\fuente{(Zac 4:7, RVR1960)}\end{cita}
\scriptsize ``No con ejército, ni con fuerza, sino con mi Espíritu'' (4:6). Dios promete ``espíritu de gracia y de oración'' (\heb{רוּחַ חֵן וְתַחֲנוּנִים}, 12:10).
\end{column}
\end{columns}
\end{frame}

\begin{frame}{``A los humildes dará gracia''}
\begin{columns}[c,onlytextwidth]
\begin{column}{0.52\textwidth}
\centering\heb{וְלַעֲנִיִּים יִתֶּן־חֵן}
\begin{cita}Ciertamente él escarnecerá a los escarnecedores, Y a los humildes dará gracia.\fuente{(Pr 3:34, RVR1960)}\end{cita}
\centering\grc{κύριος ὑπερηφάνοις ἀντιτάσσεται, ταπεινοῖς δὲ δίδωσιν χάριν}\par
{\scriptsize\color{gris}(Pr 3:34 LXX; citado en Stg 4:6 y 1 P 5:5)}
\end{column}
\begin{column}{0.46\textwidth}\small\RaggedRight
\begin{itemize}
\item La humildad no \destacar{merece} la gracia: es la \destacar{mano vacía} que la puede recibir.
\item El soberbio no la recibe porque cree que se la deben.
\end{itemize}
\end{column}
\end{columns}
\end{frame}

\begin{frame}{La gracia en el rostro de Dios}
\begin{columns}[c,onlytextwidth]
\begin{column}{0.54\textwidth}
\centering\heb{יָאֵר יְהוָה פָּנָיו אֵלֶיךָ וִיחֻנֶּךָּ}
\begin{cita}\SENOR{} te bendiga, y te guarde; \SENOR{} haga resplandecer su rostro sobre ti, y tenga de ti misericordia; \SENOR{} alce sobre ti su rostro, y ponga en ti paz.\fuente{(Nm 6:24--26, RVR1960)}\end{cita}
\end{column}
\begin{column}{0.44\textwidth}\small\RaggedRight
\begin{itemize}
\item \heb{וִיחֻנֶּךָּ} \emph{viyjunnekka}: ``y te dé \destacar{gracia}''.
\item La bendición termina en \emph{shalom}: \destacar{gracia y paz}.
\item ``El pueblo que escapó de la espada halló gracia en el desierto'' (Jer 31:2).
\end{itemize}
\end{column}
\end{columns}
\end{frame}

% ------------------------------------------------------------------
\section{Pink}

\begin{frame}{Pink: el recorrido del capítulo 13}
\centering
\begin{tikzpicture}
 \foreach \x/\y/\n/\t/\d/\r in {0/0/1/{Excluye el mérito}/{``La gracia y las obras no se unirán más que un ácido con un alcalino''}/{Ro 11:6; p. 69}, 4.1/0/2/{Tres rasgos}/{Eterna, gratuita y soberana}/{2 Ti 1:9; Ro 3:24; 5:21}, 8.2/0/3/{Distingue}/{``Gracia distintiva'': Abraham elegido de entre idólatras}/{Éx 33:19; p. 70}, 2.05/-2.45/4/{Brilla en los indignos}/{Manasés, Saulo, los corintios (Hervey)}/{1 Co 6:9--11; pp. 70--71}, 6.15/-2.45/5/{Por Cristo, por el Espíritu}/{Padre: fuente; Hijo: canal; Espíritu: sembrador}/{Jn 1:17; Zac 12:10; pp. 71--73}}{
  \node[caja,anchor=north] (b\n) at (\x,\y) {\parbox[t][1.75cm][t]{3.55cm}{\RaggedRight
    {\small\textbf{\textcolor{azul}{\t}}}\\[3pt]\scriptsize \d\par\vfill{\tiny\color{gris}(\r)}}};
  \node[circle,fill=azul,text=papiro,font=\scriptsize\bfseries,inner sep=1.6pt] at (b\n.north west) {\n};}
\end{tikzpicture}
\end{frame}

\begin{frame}{Pink: ``tan muertos''}
\begin{cita}[colback=dorado!8!papiro,colframe=dorado!85!azul]
La gracia es una provisión para aquellos hombres que están tan caídos que no pueden levantar el hacha de la justicia, tan corruptos que no pueden cambiar su propia naturaleza, tan reacios a Dios que no pueden volverse a Él\ldots{} y tan muertos que Él mismo debe abrir sus tumbas y levantarlos a la resurrección.\fuente{(G. S. Bishop, en Pink, 2020, cap. 13, p. 73)}\end{cita}
\precision{Pink escribe desde una teología reformada estricta (la gracia salvadora ``sólo sobre los elegidos''). La serie toma de él la oposición gracia/obras, que es de Pablo (Ro 11:6), y los rasgos de la gracia. Su afirmación de que ningún texto la relaciona con ``la humanidad en general'' debe leerse junto a Tit 2:11: ``la gracia de Dios se ha manifestado para salvación a todos los hombres''.}
\end{frame}

% ------------------------------------------------------------------
\section{En Israel}

\begin{frame}{Manasés: el peor rey de Judá}
\begin{columns}[t,onlytextwidth]
\begin{column}{0.56\textwidth}
\begin{cita}\ldots{}levantó altares a los baales\ldots{} Edificó también altares en la casa de \SENOR\ldots{} Y pasó sus hijos por fuego en el valle del hijo de Hinom\ldots{} se excedió en hacer lo malo ante los ojos de \SENOR, hasta encender su ira.\fuente{(2 Cr 33:3--6, RVR1960)}\end{cita}
\end{column}
\begin{column}{0.42\textwidth}\small\RaggedRight
\begin{itemize}
\item 55 años de reinado: el más largo de Judá.
\item ``Derramó\ldots{} mucha sangre inocente\ldots{} hasta llenar a Jerusalén'' (2 R 21:16).
\item ``Más mal que las naciones'' (2 Cr 33:9).
\item Gn 6:5 repetido \destacar{en Jerusalén}.
\end{itemize}
\end{column}
\end{columns}
\end{frame}

\begin{frame}{La súplica en Babilonia}
\centering\heb{וַיִּתְפַּלֵּל אֵלָיו וַיֵּעָתֶר לוֹ וַיִּשְׁמַע תְּחִנָּתוֹ}
\begin{cita}Mas luego que fue puesto en angustias, oró a \SENOR{} su Dios, humillado grandemente en la presencia del Dios de sus padres. Y habiendo orado a él, fue atendido; pues Dios oyó su oración y lo restauró a Jerusalén, a su reino. Entonces reconoció Manasés que \SENOR{} era Dios.\fuente{(2 Cr 33:12--13, RVR1960)}\end{cita}
\scriptsize\renewcommand{\arraystretch}{1.1}
\begin{tabular}{@{}>{\RaggedRight\arraybackslash}p{2.6cm}>{\RaggedRight\arraybackslash}p{9.0cm}@{}}
\heb{וַיִּכָּנַע מְאֹד} & ``se humilló mucho'': la condición de Pr 3:34. \\
\heb{תְּחִנָּתוֹ} \emph{tejinnató} & ``su pedido de \destacar{gracia}'' (RVR1960: ``su oración''). Presenta una súplica, no una obra. \\
\heb{וַיֵּדַע} & ``y \destacar{conoció}'': el verbo de Jer 9:24. La reforma (33:15--16) es \destacar{fruto}, no causa. \\
\end{tabular}\par
\precision{solo Crónicas narra el arrepentimiento (2 Cr 33:11--20); Reyes lo omite. La \emph{Oración de Manasés} es posterior y no canónica.}
\end{frame}

\begin{frame}{El escándalo: la abertura bajo el trono}
\begin{columns}[t,onlytextwidth]
\begin{column}{0.56\textwidth}
\begin{cita}Y los ángeles servidores cerraban las ventanas para que la oración de Manasés no subiera\ldots{} Les dijo: Si no lo recibo en arrepentimiento, cierro la puerta delante de todos los que se arrepienten\ldots{} Le cavó una abertura debajo de su trono de gloria y oyó su súplica.\fuente{(Talmud de Jerusalén, \emph{Sanedrín} 10:2)}\end{cita}
\centering\heb{חָתַר לוֹ חֲתִירָה מִתַּחַת כִּסֵּא הַכָּבוֹד}
\end{column}
\begin{column}{0.42\textwidth}\scriptsize\RaggedRight
\begin{itemize}
\item Misná \emph{Sanedrín} 10:2: Manasés sin parte en el mundo venidero; \destacar{Rabí Yehudá} objeta con 2 Cr 33:13.
\item Juego de letras: \heb{וַיֵּעָתֶר} (``se dejó rogar'') leído \heb{וַיֵּחָתֶר} (``le cavó''). También en \emph{Sanedrín} 103a (R. Shimón bar Yojái).
\item Manasés recordó Dt 4:30--31: ``Dios misericordioso es \SENOR{} tu Dios''.
\end{itemize}
\precision{lectura midrásica, no historia; traducción propia.}
\end{column}
\end{columns}
\end{frame}

% ------------------------------------------------------------------
\section{Judaísmo}

\begin{frame}{``Aun Noé no lo merecía''}
\begin{columns}[t,onlytextwidth]
\begin{column}{0.54\textwidth}
\centering\heb{אֲפִלּוּ נֹחַ שֶׁנִּשְׁתַּיֵּיר מֵהֶן לֹא הָיָה כְּדַאי}\par\heb{אֶלָּא שֶׁמָּצָא חֵן בְּעֵינֵי ה׳}
\begin{cita}Aun Noé, que quedó de ellos, no lo merecía, sino que halló gracia ante los ojos del \SENOR.\fuente{(\emph{Bereshit Rabá} 29:1)}\end{cita}
\end{column}
\begin{column}{0.44\textwidth}\small\RaggedRight
\begin{itemize}
\item R. Abbá bar Kahaná lee juntos ``me arrepiento de haberlos hecho, \destacar{y Noé}'' (6:7--8).
\item Escuela de R. Ismael: ``también sobre Noé se había decretado la sentencia'' (\emph{Sanedrín} 108a).
\item R. Simón: \destacar{Noé halló}; Dios no ``halló'' a Noé (BR 29:3).
\end{itemize}
\precision{otras voces del midrash atribuyen la salvación de Noé a sus méritos.}
\end{column}
\end{columns}
\end{frame}

\begin{frame}{``Don gratuito'': Rashi y el Sifré}
\centering\heb{אֵין חִנּוּן בְּכָל מָקוֹם אֶלָּא לְשׁוֹן מַתְּנַת חִנָּם}
\begin{cita}Toda forma de \emph{jinnún} significa siempre un don gratuito. Aunque los justos podrían apoyarse en sus buenas obras, no piden al Omnipresente sino un don gratuito.\fuente{(Rashi, \emph{Comentario a Dt} 3:23)}\end{cita}
\small\RaggedRight
\begin{itemize}
\item Dt 3:23: \heb{וָאֶתְחַנַּן}, ``y \destacar{supliqué gracia}''. Moisés usa la palabra que Dios usó en Éx 33:19.
\item \emph{Sifré Devarim} 26: Moisés y David ``podían sostener el mundo con sus buenas obras'' y pidieron solo \destacar{de balde} (\emph{jinnam}); ``¡cuánto más'' nosotros!
\end{itemize}
\precision{Rashi da una segunda explicación: \emph{va'etjanán} es una de las diez palabras para ``oración''.}
\end{frame}

\begin{frame}{``Aunque no lo merezca''; \emph{Avinu Malkenu}}
\begin{columns}[t,onlytextwidth]
\begin{column}{0.48\textwidth}
\centering\heb{וְחַנֹּתִי אֶת אֲשֶׁר אָחוֹן}\par\heb{אַף עַל פִּי שֶׁאֵינוֹ הָגוּן}\par\smallskip
\RaggedRight\small Rabí Meir: ``tendré gracia de quien tendré gracia, \destacar{aunque no lo merezca}'' (\emph{Berajot} 7a).\par\medskip
\scriptsize R. Akiva en la sequía: ``Padre nuestro, Rey nuestro, no tenemos rey sino a ti\ldots{} por tu causa ten compasión de nosotros'' (\emph{Taanit} 25b).
\end{column}
\begin{column}{0.50\textwidth}
\centering\heb{אָבִינוּ מַלְכֵּנוּ חָנֵּנוּ וַעֲנֵנוּ}\par\heb{כִּי אֵין בָּנוּ מַעֲשִׂים}
\begin{cita}Padre nuestro, Rey nuestro, ten gracia de nosotros y respóndenos, porque no hay en nosotros obras; haz con nosotros justicia y amor leal, y sálvanos.\fuente{(\emph{Sidur}, \emph{Avinu Malkenu})}\end{cita}
\scriptsize La oración del publicano (Lc 18:13) en la liturgia.\par
\precision{la línea \emph{jonnenu va'anenu} es parte de la forma posterior de la oración (Elbogen, 1993).}
\end{column}
\end{columns}
\end{frame}

\begin{frame}{En la oración diaria y en Qumrán}
\begin{columns}[t,onlytextwidth]
\begin{column}{0.48\textwidth}\small\RaggedRight
\textbf{\textcolor{azul}{\emph{Amidá}, 4.ª bendición}}\par
\centering\heb{אַתָּה חוֹנֵן לְאָדָם דַּעַת}\par\RaggedRight
``Tú das \destacar{por gracia} al hombre conocimiento''.\par\medskip
\textbf{\textcolor{azul}{Bendición de la mesa}}\par
\centering\heb{בְּחֵן בְּחֶסֶד וּבְרַחֲמִים}\par\RaggedRight
Dios alimenta al mundo ``con gracia, con amor leal y con compasión'': Éx 34:6 hecho \destacar{acción de gracias por el pan}.
\end{column}
\begin{column}{0.50\textwidth}
\begin{cita}En cuanto a mí, si tropiezo, las misericordias de Dios serán mi salvación eterna\ldots{} mi justificación será por la justicia de Dios\ldots{} Me acercará por su gracia, y por su misericordia traerá mi justificación.\fuente{(1QS XI, 12--14)}\end{cita}
\precision{traducción propia desde Vermes (2004). Paralelo, no dependencia, con Pablo (Dunn, 1998).}
\end{column}
\end{columns}
\end{frame}

% ------------------------------------------------------------------
\section{En Cristo}

\begin{frame}{``Justificados gratuitamente''}
\fotolado{panel_romanos3.jpg}{Elaboración propia: Ro 3:24 en escritura uncial (texto NA28).}
\begin{columns}[c,onlytextwidth]
\begin{column}{0.58\textwidth}
\begin{cita}por cuanto todos pecaron, y están destituidos de la gloria de Dios, siendo justificados gratuitamente por su gracia, mediante la redención que es en Cristo Jesús.\fuente{(Ro 3:23--24, RVR1960)}\end{cita}
\scriptsize\RaggedRight
\begin{itemize}
\item \grc{δωρεάν}: el adverbio de la LXX para \emph{jinnam} (Is 52:3).
\item \grc{τῇ αὐτοῦ χάριτι}: ``por la gracia \destacar{de él}''.
\item \grc{ἀπολύτρωσις}: rescate, como Is 52:3 (\grc{λυτρωθήσεσθε}).
\item La gloria perdida (3:23) es devuelta por gracia. Gratuita para el que recibe, \destacar{costosa} para el que da (3:25--26).
\end{itemize}
\end{column}
\begin{column}{0.40\textwidth}\end{column}
\end{columns}
\end{frame}

\begin{frame}{Abraham y el impío}
\begin{columns}[c,onlytextwidth]
\begin{column}{0.52\textwidth}
\begin{cita}Pero al que obra, no se le cuenta el salario como gracia, sino como deuda; mas al que no obra, sino cree en aquel que justifica al impío, su fe le es contada por justicia.\fuente{(Ro 4:4--5, RVR1960)}\end{cita}
\begin{cita}Por tanto, es por fe, para que sea por gracia, a fin de que la promesa sea firme para toda su descendencia\ldots\fuente{(Ro 4:16, RVR1960)}\end{cita}
\end{column}
\begin{column}{0.46\textwidth}\small\RaggedRight
\begin{itemize}
\item \grc{κατὰ χάριν ἀλλὰ κατὰ ὀφείλημα}: gracia o \destacar{deuda}.
\item Si la promesa dependiera de las obras, sería tan frágil como quien las hace; por gracia, es \destacar{tan firme como Dios}.
\item David: ``Bienaventurados aquellos cuyas iniquidades son perdonadas'' (Sal 32:1; Ro 4:7).
\end{itemize}
\end{column}
\end{columns}
\end{frame}

\begin{frame}{``Sobreabundó la gracia''}
\begin{cita}Pero la ley se introdujo para que el pecado abundase; mas cuando el pecado abundó, sobreabundó la gracia; para que así como el pecado reinó para muerte, así también la gracia reine por la justicia para vida eterna mediante Jesucristo, Señor nuestro.\fuente{(Ro 5:20--21, RVR1960)}\end{cita}
\centering
\begin{tikzpicture}[font=\scriptsize]
 \node[caja,text width=3.8cm,align=center] (a) {\textbf{Adán}\\ ``reinó la muerte'' (5:17)\\ ``el pecado reinó'' (5:21)};
 \node[draw=dorado,line width=1.6pt,fill=azul,text=papiro,rounded corners=3pt,text width=3.8cm,align=center,inner sep=5pt,right=2.4cm of a] (c) {\textbf{Cristo}\\ ``reinarán en vida'' (5:17)\\ ``la gracia reine'' (5:21)};
 \draw[-{Latex[length=3mm]},dorado,line width=1.6pt] (a) -- node[above,text=gris] {``mucho más''} (c);
\end{tikzpicture}\par\smallskip
\scriptsize \grc{ὑπερεπερίσσευσεν}: ``sobre-sobreabundó''. ``¿Perseveraremos en el pecado para que la gracia abunde? En ninguna manera'' (6:1--2): \destacar{cambio de dueño}, no permiso (6:14).
\end{frame}

\begin{frame}{Un remanente por gracia}
\begin{columns}[c,onlytextwidth]
\begin{column}{0.52\textwidth}
\begin{cita}Así también aun en este tiempo ha quedado un remanente escogido por gracia. Y si por gracia, ya no es por obras; de otra manera la gracia ya no es gracia\ldots\fuente{(Ro 11:5--6, RVR1960)}\end{cita}
\begin{cita}Por tanto, sabe que no es por tu justicia que \SENOR{} tu Dios te da esta buena tierra para tomarla; porque pueblo duro de cerviz eres tú.\fuente{(Dt 9:6, RVR1960)}\end{cita}
\end{column}
\begin{column}{0.46\textwidth}\small\RaggedRight
\begin{itemize}
\item Elías y los siete mil (1 R 19:18; Ro 11:4).
\item Pablo lee Israel \destacar{como lo lee la Torá}: elegido por amor (Dt 7:7--8), no por su justicia (Dt 9:5--6).
\end{itemize}
\precision{la ``nueva perspectiva'' (Dunn, 1998; Wright, 2013) subraya la elección por gracia en el judaísmo; Barclay (2015): todos hablaban de gracia, no todos la entendían igual.}
\end{column}
\end{columns}
\end{frame}

\begin{frame}{``Por la gracia de Dios soy lo que soy''}
\begin{columns}[c,onlytextwidth]
\begin{column}{0.54\textwidth}
\begin{cita}Porque yo soy el más pequeño de los apóstoles, que no soy digno de ser llamado apóstol, porque perseguí a la iglesia de Dios. Pero por la gracia de Dios soy lo que soy; y su gracia no ha sido en vano para conmigo, antes he trabajado más que todos ellos; pero no yo, sino la gracia de Dios conmigo.\fuente{(1 Co 15:9--10, RVR1960)}\end{cita}
\end{column}
\begin{column}{0.44\textwidth}\small\RaggedRight
\begin{itemize}
\item \grc{χάριτι δὲ θεοῦ εἰμι ὅ εἰμι}: forma \emph{idem per idem}, como Éx 3:14 y 33:19.
\item El perseguidor ``llamado por su gracia'' (Gá 1:13--15): el \destacar{Manasés} del Nuevo Testamento.
\item La gracia es \destacar{don} y \destacar{fuerza} para responder.
\end{itemize}
\precision{1 Ti 1:13--16 (autoría discutida) es eco de 1 Co 15 y Gá 1.}
\end{column}
\end{columns}
\end{frame}

\begin{frame}{Gálatas y Efesios}
\begin{columns}[t,onlytextwidth]
\begin{column}{0.48\textwidth}
\begin{cita}No desecho la gracia de Dios; pues si por la ley fuese la justicia, entonces por demás murió Cristo.\fuente{(Gá 2:21, RVR1960)}\end{cita}
\scriptsize ``Por demás'' = \grc{δωρεάν}: si la justicia se ganara, la cruz habría sido \emph{jinnam} en el peor sentido. La gracia se rechaza también \destacar{intentando pagarla}.
\end{column}
\begin{column}{0.50\textwidth}
\begin{cita}Porque por gracia sois salvos por medio de la fe; y esto no de vosotros, pues es don de Dios; no por obras, para que nadie se gloríe. Porque somos hechura suya, creados en Cristo Jesús para buenas obras\ldots\fuente{(Ef 2:8--10, RVR1960)}\end{cita}
\scriptsize Obras: \destacar{no causa} (``no por obras''), \destacar{sí propósito} (``para buenas obras''). Ef 1:6: \grc{ἐχαρίτωσεν}, ``nos hizo agraciados''.\par
\precision{autoría de Efesios discutida; su doctrina coincide con Ro y Gá.}
\end{column}
\end{columns}
\end{frame}

\begin{frame}{``Gracia y paz''}
\begin{columns}[c,onlytextwidth]
\begin{column}{0.52\textwidth}
\centering\grc{χάρις ὑμῖν καὶ εἰρήνη ἀπὸ θεοῦ πατρὸς ἡμῶν καὶ κυρίου Ἰησοῦ Χριστοῦ}
\begin{cita}Gracia y paz a vosotros, de Dios nuestro Padre y del Señor Jesucristo.\fuente{(Ro 1:7, RVR1960)}\end{cita}
\end{column}
\begin{column}{0.46\textwidth}\small\RaggedRight
\begin{itemize}
\item Saludo griego: \grc{χαίρειν} (Hch 15:23). Pablo escribe \grc{χάρις}.
\item Saludo judío: \emph{shalom}. Pablo lo une: \destacar{gracia y paz}.
\item La gracia \destacar{abre y cierra} cada carta (Ro 16:20; Gá 6:18).
\end{itemize}
\precision{eco posible de Nm 6:25--26 (``te dé gracia\ldots{} ponga en ti paz''); Pablo no cita el texto y la LXX tiene \grc{ἐλεήσαι}.}
\end{column}
\end{columns}
\end{frame}

\begin{frame}{La gracia tiene rostro}
\begin{columns}[t,onlytextwidth]
\begin{column}{0.48\textwidth}
\begin{cita}Porque ya conocéis la gracia de nuestro Señor Jesucristo, que por amor a vosotros se hizo pobre, siendo rico, para que vosotros con su pobreza fueseis enriquecidos.\fuente{(2 Co 8:9, RVR1960)}\end{cita}
\begin{cita}\ldots{}y vimos su gloria\ldots{} lleno de gracia y de verdad\ldots{} y gracia sobre gracia.\fuente{(Jn 1:14, 16, RVR1960)}\end{cita}
\end{column}
\begin{column}{0.50\textwidth}
\scriptsize\renewcommand{\arraystretch}{1.15}
\begin{tabular}{@{}>{\RaggedRight\arraybackslash}p{1.3cm}>{\RaggedRight\arraybackslash}p{4.4cm}@{}}
\toprule
Éx 34:6 & \heb{וְרַב־חֶסֶד וֶאֱמֶת} \\
LXX & \grc{πολυέλεος καὶ ἀληθινός} \\
Jn 1:14 & \grc{πλήρης χάριτος καὶ ἀληθείας} \\
\bottomrule
\end{tabular}\par\medskip
\small\RaggedRight Moisés pidió ver la gloria (Éx 33:18); Juan dice que se vio en el Verbo. Donde la LXX tenía \grc{ἔλεος}, Juan escribe \grc{χάρις}.\par
\precision{eco reconocido, no cita. Juan no es Pablo; se incluye porque Pink se apoya en Jn 1:17.}
\end{column}
\end{columns}
\end{frame}

\begin{frame}{``Bástate mi gracia''; la gracia que enseña}
\begin{columns}[t,onlytextwidth]
\begin{column}{0.48\textwidth}
\begin{cita}Y me ha dicho: Bástate mi gracia; porque mi poder se perfecciona en la debilidad\ldots\fuente{(2 Co 12:9, RVR1960)}\end{cita}
\scriptsize \grc{Ἀρκεῖ σοι ἡ χάρις μου}: la gracia que justifica al impío \destacar{sostiene} al débil. Pablo se gloría solo en lo que muestra el poder de Dios (Jer 9:24).
\end{column}
\begin{column}{0.50\textwidth}
\begin{cita}Porque la gracia de Dios se ha manifestado para salvación a todos los hombres, enseñándonos que, renunciando a la impiedad y a los deseos mundanos, vivamos en este siglo sobria, justa y piadosamente\ldots\fuente{(Tit 2:11--12, RVR1960)}\end{cita}
\scriptsize \grc{ἐπεφάνη}: ``apareció'', como una luz. \grc{παιδεύουσα}: ``educando''. La gracia que perdona \destacar{también forma}.
\end{column}
\end{columns}
\end{frame}

% ------------------------------------------------------------------
\section{Vivirla}

\begin{frame}{``De gracia recibisteis, dad de gracia''}
\begin{columns}[t,onlytextwidth]
\begin{column}{0.48\textwidth}
\begin{cita}\ldots{}de gracia recibisteis, dad de gracia.\fuente{(Mt 10:8, RVR1960)}\end{cita}
\scriptsize \grc{δωρεὰν ἐλάβετε, δωρεὰν δότε}. Pablo predicó \grc{δωρεάν} (2 Co 11:7) y llamó ``gracia'' a la ofrenda para los pobres (2 Co 8:1, 7).
\begin{cita}Sea vuestra palabra siempre con gracia, sazonada con sal\ldots\fuente{(Col 4:6, RVR1960)}\end{cita}
\end{column}
\begin{column}{0.50\textwidth}
\begin{cita}Antes sed benignos unos con otros, misericordiosos, perdonándoos unos a otros, como Dios también os perdonó a vosotros en Cristo.\fuente{(Ef 4:32, RVR1960)}\end{cita}
\scriptsize ``Perdonar'' = \grc{χαρίζομαι}: \destacar{dar gracia}. ``Si amáis a los que os aman, ¿qué mérito (\grc{χάρις}) tenéis?'' (Lc 6:32): la gracia verdadera se da a quien no la devuelve.
\end{column}
\end{columns}
\end{frame}

\begin{frame}{Cinco respuestas}
\centering
\begin{tikzpicture}
 \foreach \x/\y/\n/\t/\d/\r in {0/0/1/{Manos vacías}/{Recibir como Manasés: con una \emph{tejinná}, no con un pago}/{Ro 4:4--5; Is 55:1}, 4.1/0/2/{No es permiso}/{La gracia que perdona enseña a renunciar}/{Ro 6:1--2; Tit 2:12}, 8.2/0/3/{Gratitud}/{``¿Qué tienes que no hayas recibido?''}/{1 Co 4:7; Stg 4:6}, 2.05/-2.35/4/{Dar y perdonar}/{La medida es lo recibido, no lo merecido}/{Mt 10:8; Ef 4:32}, 6.15/-2.35/5/{El trono de la gracia}/{``hallar gracia para el oportuno socorro''}/{He 4:16}}{
  \node[caja,anchor=north] (b\n) at (\x,\y) {\parbox[t][1.65cm][t]{3.55cm}{\RaggedRight
    {\small\textbf{\textcolor{azul}{\t}}}\\[3pt]\scriptsize \d\par\vfill{\tiny\color{gris}(\r)}}};
  \node[circle,fill=azul,text=papiro,font=\scriptsize\bfseries,inner sep=1.6pt] at (b\n.north west) {\n};}
\end{tikzpicture}
\end{frame}

\begin{frame}{Para la reflexión}
\begin{cita}Acerquémonos, pues, confiadamente al trono de la gracia, para alcanzar misericordia y hallar gracia para el oportuno socorro.\fuente{(He 4:16, RVR1960)}\end{cita}
\small
\begin{itemize}
\item ¿Hay algo que todavía intento \destacar{pagarle} a Dios para sentirme aceptado?
\item ¿Me escandaliza, como a Jonás, la gracia que Dios muestra a otros?
\item ¿A quién debo hoy un perdón ``de balde'', sin esperar que lo devuelva?
\item ¿Mis palabras dan gracia a los que las oyen (Ef 4:29)?
\end{itemize}
\end{frame}

% ------------------------------------------------------------------
\section{Síntesis}

\begin{frame}{Síntesis}
\scriptsize\renewcommand{\arraystretch}{1.1}
\begin{tabular}{@{}>{\RaggedRight\arraybackslash}p{1.8cm}>{\RaggedRight\arraybackslash}p{6.6cm}>{\RaggedRight\arraybackslash}p{3.2cm}@{}}
\toprule
\textbf{Lente} & \textbf{Aporte} & \textbf{Textos} \\ \midrule
Texto ancla & la gracia aparece en el peor momento de la humanidad & Gn 6:5--8; 8:21 \\
La palabra & \emph{jen}: favor en los ojos del superior; \emph{jinnam}: de balde & Rt 2:10; Éx 21:2 \\
El Nombre & \emph{jannun}: oye al que no puede reclamar & Éx 22:27; 34:6 \\
Pink & excluye el mérito; eterna, gratuita, soberana & Ro 11:6; 5:21 \\
Israel & Manasés oído en su súplica & 2 Cr 33:12--13 \\
Judaísmo & ``don gratuito''; ``no hay en nosotros obras'' & Dt 3:23; \emph{Avinu Malkenu} \\
Pablo & justificados gratuitamente; la gracia reina & Ro 3:24; 5:21; 1 Co 15:10 \\
Vida & dar, perdonar y hablar con gracia & Mt 10:8; Ef 4:32 \\
\bottomrule
\end{tabular}\par\medskip
\begin{tcolorbox}[colback=papiro,colframe=dorado,boxrule=1pt,arc=3pt,fontupper=\small\itshape]
La gloria que el pecador perdió (Ro 3:23) le es dada otra vez por gracia: ``gracia y gloria dará \SENOR'' (Sal 84:11). Esa plenitud habita en Cristo (Col 2:9).
\end{tcolorbox}
\end{frame}

% ------------------------------------------------------------------
\section{Glosario}

\begin{frame}{Glosario}
\scriptsize
\begin{columns}[t,onlytextwidth]
\begin{column}{0.49\textwidth}
\begin{description}\setlength{\itemsep}{1pt}
\item[\emph{Avinu Malkenu}] (\heb{אָבִינוּ מַלְכֵּנוּ}) ``Padre nuestro, Rey nuestro''.
\item[\emph{Dōreán}] (\grc{δωρεάν}) gratuitamente, de balde (Ro 3:24).
\item[\emph{Jannun}] (\heb{חַנּוּן}) clemente (Éx 22:27; 34:6).
\item[\emph{Janán}] (\heb{חָנַן}) mostrar gracia; hitpael: suplicar (Dt 3:23).
\item[\emph{Járis}] (\grc{χάρις}) gracia, favor, agradecimiento.
\item[\emph{Járisma}] (\grc{χάρισμα}) don, dádiva (Ro 6:23).
\end{description}
\end{column}
\begin{column}{0.49\textwidth}
\begin{description}\setlength{\itemsep}{1pt}
\item[\emph{Jen}] (\heb{חֵן}) gracia, favor, encanto (Gn 6:8).
\item[\emph{Jinnam}] (\heb{חִנָּם}) de balde, gratis, sin causa (Is 52:3).
\item[\emph{Matnat jinnam}] (\heb{מַתְּנַת חִנָּם}) don gratuito (Rashi).
\item[\emph{Opsōnia}] (\grc{ὀψώνια}) paga del soldado (Ro 6:23).
\item[\emph{Tejinná}] (\heb{תְּחִנָּה}) súplica, pedido de gracia (2 Cr 33:13).
\end{description}
\end{column}
\end{columns}
\end{frame}

% ------------------------------------------------------------------
\section{Referencias}

\begin{frame}{Referencias}
\fontsize{5.0}{6.0}\selectfont
\begin{columns}[t,onlytextwidth]
\begin{column}{0.49\textwidth}
\refapa{Barclay, J.~M.~G. (2015). \emph{Paul and the gift}. Eerdmans.}
\refapa{Bauer, W., Danker, F.~W., Arndt, W.~F., \& Gingrich, F.~W. (2000). \emph{A Greek-English lexicon of the New Testament} (3.ª ed.). University of Chicago Press. [BDAG]}
\refapa{\emph{Bereshit Rabá} [29:1, 3]. (s.f.). Sefaria.}
\refapa{Berkhof, L. (1938). \emph{Systematic theology}. Eerdmans.}
\refapa{Brown, F., Driver, S.~R., \& Briggs, C.~A. (1906). \emph{A Hebrew and English lexicon of the Old Testament}. Clarendon Press.}
\refapa{Cassuto, U. (1961). \emph{A commentary on the book of Genesis: Part I}. Magnes Press.}
\refapa{Dunn, J.~D.~G. (1988). \emph{Romans 1--8} (WBC 38A). Word Books.}
\refapa{Dunn, J.~D.~G. (1998). \emph{The theology of Paul the apostle}. Eerdmans.}
\refapa{Elbogen, I. (1993). \emph{Jewish liturgy: A comprehensive history}. Jewish Publication Society.}
\refapa{Elliger, K., \& Rudolph, W. (Eds.). (1997). \emph{Biblia Hebraica Stuttgartensia} (5.ª ed.). Deutsche Bibelgesellschaft.}
\refapa{Guggenheimer, H.~W. (Ed. y Trad.). (1999--2015). \emph{The Jerusalem Talmud} [Sanedrín 10:2]. De Gruyter. Sefaria.}
\refapa{Koehler, L., \& Baumgartner, W. (2001). \emph{The Hebrew and Aramaic lexicon of the Old Testament}. Brill. [HALOT]}
\refapa{Moo, D.~J. (1996). \emph{The epistle to the Romans}. Eerdmans.}
\end{column}
\begin{column}{0.49\textwidth}
\refapa{Nestle, E., \& Aland, K. (Eds.). (2012). \emph{Novum Testamentum Graece} (28.ª ed.). Deutsche Bibelgesellschaft.}
\refapa{Pink, A.~W. (2020). \emph{Los atributos de Dios} (E. Rodríguez Cruz, Trad.). Chapel Library. (Obra original publicada en 1930)}
\refapa{Rahlfs, A., \& Hanhart, R. (Eds.). (2006). \emph{Septuaginta} (Ed. rev.). Deutsche Bibelgesellschaft.}
\refapa{Rashi. (s.f.). \emph{Comentario a la Torá} [Dt 3:23]. Sefaria.}
\refapa{\emph{Santa Biblia: Reina-Valera 1960}. (1960). Sociedades Bíblicas en América Latina.}
\refapa{Sasson, J.~M. (1975). Word-play in Gen 6:8--9. \emph{The Catholic Biblical Quarterly, 37}(2), 165--166.}
\refapa{\emph{Shemot Rabá} [3:6]; \emph{Sifré Devarim} [26]; Misná [Sanedrín 10:2]. (s.f.). Sefaria.}
\refapa{\emph{Sidur Ashkenaz}; \emph{Sidur Sefard} [Avinu Malkenu]. (s.f.). Sefaria.}
\refapa{Steinsaltz, A. (Ed.). (2017). \emph{Koren Talmud Bavli: Edición William Davidson} [Berajot 7a; Taanit 25b; Sanedrín 103a; 108a]. Sefaria.}
\refapa{Vermes, G. (2004). \emph{The complete Dead Sea Scrolls in English} (Ed. rev.). Penguin.}
\refapa{Wright, N.~T. (2013). \emph{Paul and the faithfulness of God}. Fortress Press.}
\par\smallskip{\color{gris}Citas bíblicas: RVR1960; ``Jehová'' se reemplaza por \SENOR. Hebreo: BHS; griego: NA28 y LXX. Lo rabínico y las lecturas propias se marcan como \emph{precisión}.}
\end{column}
\end{columns}
\end{frame}

\end{document}
